\fig{fv:k3:rope}{活动滑轮的两段绳同时改变长度}{\begin{tikzpicture}[scale=.9]
\draw[thick](-1,2.5)--(4.5,2.5);\draw[thick](0,2.5)--(0,0)arc(180:360:1)--(2,2.5);
\draw[thick](1,0)circle(1);\draw[thick](1,0)--(1,-1.5);\draw[thick](2,2.5)--(3,2.5)--(3,1.2);
\draw[->,blue](3.5,2.1)--(3.5,1.1)node[midway,right]{$v_f$};\draw[->,red](1.6,-1)--(1.6,-1.8)node[right]{$v_c$};
\draw[<->](-.5,2.5)--(-.5,0)node[midway,left]{$y$};\node[below]at(3,.8){自由端};
\end{tikzpicture}}
图\ref{fv:k3:rope}按理想定包角模型画出两条支承活动滑轮的直段。若从固定横梁向下量轮心位移 $y$和自由端位置 $z$，总绳长的变量部分为 $2y+z$，故 $2\dot y+\dot z=0$。图上的速度箭头只表示坐标的正向，求得的符号再决定真实运动方向。活动滑轮不是只改变方向的固定滑轮：轮心每下移一段，两侧各增加同样的绳长，倍数二正由这两个长度变化产生。
